% Draft Methods and Results for the DO-like Asian monsoon transition paper.
% Citations are author-year placeholders for later bibliography integration.

\documentclass[11pt]{article}
\usepackage[margin=1in]{geometry}
\usepackage{amsmath}
\usepackage[version=4]{mhchem}
\usepackage[T1]{fontenc}
\usepackage{lmodern}

\title{Methods and Results Draft}
\author{}
\date{}

\begin{document}
\maketitle

\section*{Methods}

We analyzed the weak and strong monsoon-start chronologies identified by Rousseau et al. (2023) from the Chinese speleothem composite of Cheng et al. (2016). The two catalogues were treated as separate event sequences because they represent opposite abrupt transitions in the Asian monsoon record rather than paired observations from a single two-state model. The analysis was restricted to 0--640 ka BP, giving 96 strong monsoon starts and 94 weak monsoon starts. All tests below were applied independently to the two event catalogues.

Before fitting the main event-hazard models, we performed two supporting event-timing tests. First, we used a Lohmann-style stationary occurrence test to ask whether the weak and strong start sequences are compatible with stationary random timing. Second, we used circular Rayleigh tests to ask whether event phases are uniformly distributed with respect to precession and obliquity. These tests provide diagnostic context for the hazard analysis, but they do not adjust for glacial background state; full details are given in Text S1 and Text S2 of the Supporting Information.

The main analysis used a binned Poisson hazard generalized linear model to test whether orbital phase explains transition probability after accounting for glacial background state. The 0--640 ka interval was divided into bins of width \(\Delta t = 0.2\) kyr. For each event catalogue and bin \(i\), the response variable was the event count \(Y_i\). Conditional on the covariates, event counts were modeled as independent Poisson variables:
\[
Y_i \sim \mathrm{Poisson}(\mu_i),
\qquad
\mu_i = \Delta t_i \lambda_i,
\]
where \(\lambda_i\) is the fitted event rate per kyr. The log link was applied to the rate rather than directly to the expected bin count:
\[
\log \lambda_i
= \beta_0 + \mathbf{x}_i^\mathsf{T}\boldsymbol{\beta}.
\]
Equivalently,
\[
\log \mu_i
= \log \Delta t_i + \beta_0 + \mathbf{x}_i^\mathsf{T}\boldsymbol{\beta},
\]
so \(\log\Delta t_i\) enters the model as an offset. The Poisson log-likelihood was
\[
\ell(\boldsymbol{\beta})
=
\sum_i
\left[
Y_i\log\mu_i - \mu_i - \log(Y_i!)
\right].
\]

The baseline model set consisted of four candidate models: a stationary intercept-only model, a climate-background model with LR04 and atmospheric \(\ce{CO2}\), a precession-phase-only model, and a combined model:
\[
\log \lambda_i
= \beta_0
+ \beta_1 \mathrm{LR04}_i
+ \beta_2 \ce{CO2}_i
+ \beta_3 \sin\theta_i
+ \beta_4 \cos\theta_i,
\]
where \(\theta_i\) is the precession phase at the bin center. Precession phase was defined from the precession-index waveform by assigning phase \(0\) to local minima and phase \(\pi\) to local maxima, with linear interpolation of unwrapped phase between adjacent extrema. LR04 and \(\ce{CO2}\) were interpolated to bin centers and range-scaled to have zero mean and unit observed range. Precession phase was represented by \(\sin\theta_i\) and \(\cos\theta_i\), because phase is circular and cannot be treated as a linear predictor.

Models with predictors were fitted by maximum likelihood using L-BFGS-B optimization; the stationary model has the analytic maximum-likelihood rate \(N/T\). Models were ranked using AICc. Nested models were compared with likelihood-ratio tests, using
\[
\mathrm{LR} = 2(\ell_\mathrm{full}-\ell_\mathrm{reduced})
\]
and a chi-square reference distribution with degrees of freedom equal to the number of added coefficients. The key comparison was ``precession phase after climate'', in which the reduced model was LR04 + \(\ce{CO2}\) and the full model added the two circular precession-phase terms.

For models containing \(\sin\theta\) and \(\cos\theta\), we converted the two phase coefficients into an interpretable preferred phase. If the phase component is
\[
\beta_s\sin\theta+\beta_c\cos\theta,
\]
then the phase of maximum fitted rate is
\[
\theta^\ast = \operatorname{atan2}(\beta_s,\beta_c)\ \mathrm{mod}\ 2\pi,
\]
and the phase-response amplitude is
\[
A = (\beta_s^2+\beta_c^2)^{1/2}.
\]
The fitted maximum-to-minimum rate ratio over one precession cycle is therefore \(\exp(2A)\).

We performed two sensitivity experiments for the hazard-model results. First, we asked whether other plausible forcings add explanatory power beyond the baseline LR04 + \(\ce{CO2}\) + precession-phase model. The extra predictors were Antarctic temperature (AT), obliquity, eccentricity, and \(65^\circ\)N summer-solstice insolation. Each extra predictor was range-scaled in the same way as LR04 and \(\ce{CO2}\). We tested the contribution of the extra forcings in three ways: adding each extra predictor to the baseline model one at a time, adding all four extra predictors jointly, and dropping each extra predictor from the full extended model to estimate its unique contribution after the other extra predictors were present.

Second, we tested sensitivity to the bin width used in the binned hazard approximation. The baseline bin width was 0.2 kyr, chosen to keep bins short relative to expected changes in the underlying hazard. To check that this choice did not determine the conclusions, we repeated the core model comparison for bin widths of 0.2, 0.4, 0.6, 0.8, and 1.0 kyr. For each bin width, we tracked model ranking by AICc, the likelihood-ratio \(p\) value for precession phase after LR04 + \(\ce{CO2}\), the preferred precession phase, and the sparsity of the event-count bins.

\section*{Results}

The supporting event-timing tests motivate a non-stationary, phase-aware hazard model. The Lohmann-style test shows that the weak monsoon-start sequence is inconsistent with stationary random occurrence under both the unconditional Poisson and fixed-\(N\) nulls, whereas the strong monsoon-start sequence is significant under the fixed-\(N\) null and marginal under the unconditional Poisson null. Rayleigh tests further show significant precession-phase clustering for both strong and weak starts, but no significant obliquity-phase clustering. Detailed statistics and test definitions are reported in Text S1 and Text S2 of the Supporting Information.

The binned Poisson hazard models confirm that precession phase contributes information beyond glacial background state. For both weak and strong monsoon starts, the combined LR04 + \(\ce{CO2}\) + precession-phase model had the lowest AICc among the core candidate models. For strong monsoon starts, adding LR04 + \(\ce{CO2}\) to the stationary model improved the likelihood (\(p=0.0092\)), and adding precession phase to LR04 + \(\ce{CO2}\) improved it further (\(\mathrm{LR}=15.10\), \(p=5.25\times10^{-4}\), \(\Delta\mathrm{AICc}=-11.09\)). The fitted preferred precession phase in the full model was \(10.8^\circ\), close to the Rayleigh mean direction, and the fitted maximum-to-minimum rate ratio over the precession cycle was 3.18.

Weak monsoon starts show the same model-ranking pattern but with the opposite preferred precession sector. The LR04 + \(\ce{CO2}\) model improved on stationarity (\(p=4.31\times10^{-4}\)), and precession phase still improved the model after LR04 + \(\ce{CO2}\) were included (\(\mathrm{LR}=10.66\), \(p=0.00485\), \(\Delta\mathrm{AICc}=-6.65\)). The full model preferred a precession phase of \(237.5^\circ\), again close to the Rayleigh mean direction for weak starts, with a fitted maximum-to-minimum rate ratio of 2.63. These results indicate that precession phase is not merely reproducing the slow LR04/\(\ce{CO2}\) background structure; it remains associated with event probability after those background variables are included.

The extra-forcing sensitivity experiment does not identify an additional independent driver comparable to precession phase. After the baseline LR04 + \(\ce{CO2}\) + precession-phase model was fitted, adding AT, obliquity, eccentricity, or \(65^\circ\)N summer-solstice insolation one at a time did not significantly improve either event model. For strong monsoon starts, the single-extra \(p\) values ranged from 0.227 to 0.936; for weak monsoon starts, they ranged from 0.363 to 0.689. Adding all four extra forcings jointly also failed to improve the baseline model for strong starts (\(p=0.682\), \(\Delta\mathrm{AICc}=+5.75\)) or weak starts (\(p=0.733\), \(\Delta\mathrm{AICc}=+6.02\)). Drop-one tests within the full extended model likewise showed no significant unique contribution from any extra forcing. Thus, in this event-hazard framework, the robust orbital signal is the precession phase term already present in the baseline model, rather than obliquity, eccentricity, Antarctic temperature, or \(65^\circ\)N summer-solstice insolation as independent additions.

The bin-width sensitivity experiment shows that the 0.2 kyr bin choice is not driving the qualitative conclusions. Across bin widths from 0.2 to 1.0 kyr, the LR04 + \(\ce{CO2}\) + precession-phase model remained the best core model by AICc for both event catalogues. The likelihood-ratio test for precession phase after LR04 + \(\ce{CO2}\) stayed significant at every tested bin width. For strong monsoon starts, the preferred phase stayed between \(10.5^\circ\) and \(11.5^\circ\), with rate ratios of approximately 3.15--3.19. For weak monsoon starts, the preferred phase stayed between \(236.6^\circ\) and \(238.4^\circ\), with rate ratios of approximately 2.63--2.69. The preferred phases shifted by less than \(1^\circ\) relative to the 0.2 kyr baseline. The 0.2 kyr bins are sparse, with about 97\% empty bins and no multi-event bins, but wider bins also preserve the same phase preference and model ranking. We therefore treat 0.2 kyr as a conservative baseline that keeps the within-bin constant-hazard approximation short while yielding conclusions that are stable across coarser binning choices.

\end{document}
